\documentclass[a4paper,11pt]{extarticle}


\usepackage[T1]{fontenc}
\usepackage[utf8]{inputenc}
\usepackage{lmodern}

\usepackage[]{parskip}
\usepackage{hyperref, wrapfig, float, amsmath, adjustbox, array, pdflscape, soul, color, extsizes}
\usepackage{booktabs}
\usepackage{lscape}
\usepackage{setspace}
\usepackage{multirow}

\renewcommand{\arraystretch}{2}
\renewcommand{\contentsname}{Spis treści}
\renewcommand{\refname}{Bibliografia}
\renewcommand{\figurename}{Rys.}
\renewcommand{\tablename}{Tab.}

\usepackage{geometry}
\geometry{
  a4paper,
  margin=30mm,
  bottom=30mm
}

\usepackage{graphicx}
\usepackage[font=footnotesize,labelfont=bf,justification=centering]{caption}
\usepackage{subcaption}
\graphicspath{ {./images/}{../images/} }

\newcommand{\mathcolorbox}[1]{
  \begin{center}
    \colorbox{yellow}{\(\displaystyle #1\)}
  \end{center}
}
\newcommand{\pvec}[1]{
  \vec{#1}\mkern2mu\vphantom{#1}
}
\newcommand{\angstrom}{\mbox{\normalfont\AA}}

\begin{document}
  \begin{titlepage}
    \begin{center}
      \begin{figure}[H]
        \centering
        
\includegraphics[scale=.125]{logopp.png}
      \end{figure}
      \vspace*{.05\textheight}
      \huge\textbf{Politechnika Poznańska}\\
      \large\textbf{Wydział Inżynierii Materiałowej i Fizyki Technicznej}\\
      \large\textbf{Fizyka Techniczna}\\
      \vspace*{.2\textheight}
      \huge\textbf{Oznaczenie czasu połowicznego rozpadu dla naturalnego nuklidu \( ^{40}K \)}\\
      \vspace*{.025\textheight}
      \large\textbf{Jakub Bąbelek\\jbbabelek@gmail.com\\(nr.\ albumu: 136454)}
    \end{center}
  \end{titlepage}

  \doublespacing

  %\tableofcontents
  
  \clearpage

  \section{Wstęp}
  Nuklid to jądro atomowe o określonej liczbie nukleonów oraz stanie energii
  jądrowej. Liczbę protonów wyraża liczba atomowa \( Z \), natomiast liczbę
  neutronów \( N \) możemy wyznaczyć ze wzoru na liczbę masową \( A = Z + N \).

  \subsection{Rozpad promieniotwórczy}
  Rozpadem jądra promieniotwórczego określamy jego przemianę w nuklid posiadający
  niższy stan energetyczny. Wypromieniowuje on przy tym liczne cząsteczki które
  mogą indukować kolejne rozszczepienia. Podstawowym parametrem opisującym rozpad
  atomowy jest stała rozpadu \( \lambda \) jest ona miarą prawdopodobieństwa jego
  wystąpienia w jednostce czasu. Aktywność preparatu promieniotwórczego określamy
  wzorem (\ref{eq:act})

  \begin{equation}
    \alpha = - \dfrac{dN}{dt} = \lambda N
    \label{eq:act}
  \end{equation}
  \begin{center}
    \small{\( gdzie:\ N\ -\ liczba\ jąder\ atomowych,\ t\ -\ czas,\ \lambda\ -\ stała\ rozpadu \)}
  \end{center}

  Możemy również wyznaczyć czas połowicznego rozpadu tj. czas po którym liczba
  radioaktywnych jąder zmniejszy się o połowę. Znając tą zależność możemy określić
  czas połowicznego rozpadu (\ref{eq:time})

  \begin{equation}
    \dfrac{N}{N_0} = \dfrac{1}{2} = e^{-\lambda T} \Rightarrow T = \dfrac{ln(2)}{\lambda} \approx \dfrac{0,693}{\lambda}
    \label{eq:time}
  \end{equation}

  \section{Dane i obliczenia}

  \begin{table}[H]
    \centering
    \begin{adjustbox}{max width = \textwidth}
      \begin{tabular}{cccc}
      \toprule
      \textbf{masa \( [g] \)} & \textbf{impulsy \( [\dfrac{1}{s}] \)} & \textbf{srednio imp. \( [\dfrac{1}{s}] \)} & \textbf{aktywność bez tła \( [\dfrac{1}{s}] \)} \\
      \midrule
      \multirow{3}{*}{0}      & 156                   & \multirow{3}{*}{160,3}    & \multirow{3}{*}{0}              \\
                              & 157                   &                           &                                 \\
                              & 168                   &                           &                                 \\
      \midrule
      \multirow{3}{*}{0,005}  & 165                   & \multirow{3}{*}{170,7}    & \multirow{3}{*}{10,3}           \\
                              & 178                   &                           &                                 \\
                              & 169                   &                           &                                 \\
      \midrule
      \multirow{3}{*}{0,008}  & 187                   & \multirow{3}{*}{170,7}    & \multirow{3}{*}{10,3}           \\
                              & 161                   &                           &                                 \\
                              & 164                   &                           &                                 \\
      \midrule
      \multirow{3}{*}{0,015}  & 161                   & \multirow{3}{*}{180,7}    & \multirow{3}{*}{20,3}           \\
                              & 197                   &                           &                                 \\
                              & 184                   &                           &                                 \\
      \midrule
      \multirow{3}{*}{0,01}   & 227                   & \multirow{3}{*}{211,3}    & \multirow{3}{*}{51}             \\
                              & 191                   &                           &                                 \\
                              & 216                   &                           &                                 \\
      \bottomrule
      \end{tabular}
    \end{adjustbox}
  \end{table}



  \begin{figure}[H]
    \centering
    \includegraphics[width=15cm]{wykres.png}
    \caption{Zależnosc aktywności bez tła od masy KCl}
    \label{fig:wykres}
  \end{figure}

  Wartość współczynnika nachylenia prostej dla regresji liniowej przecinającej
  wykres w punkcie \( (0,0) \) z pominięciem pomiaru dla masy \( 0,01 [g] \) wynosi:
  \( a = 1396,2 \left[ \dfrac{imp}{100s * g} \right] \)

  \begin{equation}
    \alpha_{KCl} = \dfrac{1396,2}{43\%} = 3247 \left[ \dfrac{imp}{100s} \right] = 32,5 \left[ \dfrac{imp}{s} \right]
    \label{eq:1}
  \end{equation}

  \begin{equation}
    m_K = \dfrac{39,1 * 1 [g]}{74,55} = 0,52 [g]
    \label{eq:2}
  \end{equation}

  \begin{equation}
    \alpha_K = \dfrac{32,5}{0,52} = 62,5 \left[ \dfrac{imp}{s} \right]
    \label{eq:3}
  \end{equation}

  \begin{equation}
    \alpha_{^{40}K} = \dfrac{62,5}{0,00012} = 517083 \left[ \dfrac{imp}{s} \right]
    \label{eq:4}
  \end{equation}

  \begin{equation}
    \lambda = \dfrac{\alpha}{N} = \dfrac{517083}{6,02 * 10^{23}} = 8,59 * 10^{-19} \left[ \dfrac{1}{s} \right]
    \label{eq:5}
  \end{equation}

  \begin{equation}
    T = \dfrac{ln 2}{\lambda} \approx \dfrac{0,693}{8,59 * 10^{-19}} \approx 8,07 * 10^{17} [s] \approx 2,56 * 10^{10} [lat]
    \label{eq:wynik}
  \end{equation}

  Obliczamy liczbę cząstek \( \beta \) z uwzględnieniem sprawności licznika
  (\ref{eq:1}). Następnie znając masę cząsteczkową KCl \( 74,55u \) oraz masę
  atomową potasu obliczamy masę potasu w jednym gramie KCl (\ref{eq:2}).
  Znając obie wartości możemy obliczyć ilość cząstek \( \beta \) przypadającą na
  gram potasu (\ref{eq:3}). Uwzględniając fakt, że nuklid \( ^{40}K \) znajduje
  się w naturalnym potasie w bardzo małej ilości wynoszącej \( 0,012\% \)
  (\ref{eq:4}) podstawiamy uzyskaną wartość do wzoru (\ref{eq:act}) uzyskując
  stałą rozpadu (\ref{eq:5}). Końcowy wynik uzyskujemy przy użyciu wzoru (\ref{eq:time})
  który daje nam czas połowicznego rozpadu nuklidu \( ^{40}K \) (\ref{eq:wynik}).

  \section{Wnioski}
  Z obliczenia (\ref{eq:wynik}) wynika, że czas połowicznego rozpadu \( ^{40}K \)
  wynosi \( 2,56 * 10^{10} \) lat. Jest to cały rząd różnicy od wartości tablicowej
  wynoszącej \( 1,3 * 10^9 \). Dokladność wyniku może być niska ze względu na małą
  próbę, wykonanie większej ilości pomiarów mogłoby zbliżyc wynik do wartości
  tablicowej. Ponadto rozbieżność wyników mogła być spowodowana błędem ludzkim
  podczas przeprowadzania doświadczenia czy słabą dokładnością urządzeń pomiarowych.

\end{document}